\documentclass[11pt,a4paper]{article}
\usepackage[T1]{fontenc}
\usepackage[utf8]{inputenc}
\usepackage{amsmath,amssymb,amsthm}
\usepackage[margin=1in]{geometry}
\usepackage[colorlinks=true,linkcolor=blue,urlcolor=blue]{hyperref}
\theoremstyle{plain}
\newtheorem{theorem}{Theorem}
\newtheorem{lemma}[theorem]{Lemma}
\newtheorem{proposition}[theorem]{Proposition}
\newtheorem{corollary}[theorem]{Corollary}
\theoremstyle{definition}
\newtheorem{remark}[theorem]{Remark}
% A quoted conjecture can be a single hundred-term formula whose terms are thirty-digit
% numbers. TeX cannot hyphenate those, so without a little stretch every such quote runs off
% the page; the linked-a-file papers are all of that shape.
\setlength{\emergencystretch}{3em}

\title{The empirical recurrence for OEIS A264418, proved by transfer matrix}
\author{Adrian Perez Fontelles\\ \small Independent researcher}
\date{14 September 2026}
\begin{document}
\maketitle

\begin{abstract}
OEIS A264418 counts the arrangements of $0,1,\dots,LW-1$ in a grid of $W=4$ rows in which
every value moves from its home cell by an index change on a short list, and it carries an
empirical recurrence of order $70$ contributed by R.~H.~Hardin. It is true, and it is
decidable rather than empirical. Such an arrangement is a perfect matching between cells and
values, and because a value's row moves by at most $R=2$, the matching splits into
independent decisions row by row: after the cells of rows $0,\dots,i$ are filled,
every value whose home row is at most $i-R$ has been used, and only $2R$ home rows are
partly used. Carrying those as bitmasks makes the count a walk on $S=65536$ vertices, so it is
$C$-finite and the conjectured recurrence is decided by a finite exact computation.
\end{abstract}

\noindent\small 2020 Mathematics Subject Classification. 05A05, 05A15, 15A15.\normalsize

\section{The conjecture}

OEIS A264418 is ``Number of (n+1)X(3+1) arrays of permutations of 0..n*4+3 with each element having directed index change 0,1 2,2 1,0 -1,2 -2,-1 or -1,-1.''. It has offset $1$ and begins
\[
0, 12, 96, 529, 4843, 44264, 352713, 2840274,\ \dots
\]
The entry states, as a conjecture contributed by its author and never marked settled:
\begin{quote}
Empirical recurrence of order 70 (see link above)

Empirical: a(n) = 7*a(n-1) - 3*a(n-2) + 82*a(n-3) + 283*a(n-4) - 1154*a(n-5) + 4042*a(n-6) + 17289*a(n-7) - 22185*a(n-8) + 32767*a(n-9) + 188026*a(n-10) + 232472*a(n-11) - 37483*a(n-12) + 1140203*a(n-13) - 150763*a(n-14) + 7561993*a(n-15) - 22181996*a(n-16) - 58336867*a(n-17) - 54942280*a(n-18) - 61937490*a(n-19) + 338713215*a(n-20) - 616216060*a(n-21) - 750478274*a(n-22) + 3004008536*a(n-23) + 1358005083*a(n-24) - 3290415548*a(n-25) - 5332099311*a(n-26) + 18835203759*a(n-27) - 7620056042*a(n-28) + 24271726336*a(n-29) - 21686365629*a(n-30) + 23070567658*a(n-31) - 557733773*a(n-32) + 145268540287*a(n-33) + 48269697816*a(n-34) + 145549126524*a(n-35) + 209871190816*a(n-36) + 171564703835*a(n-37) + 322482128177*a(n-38) + 392962918650*a(n-39) + 322521808333*a(n-40) + 509779938296*a(n-41) - 612220254822*a(n-42) + 615594893007*a(n-43) - 1246199796257*a(n-44) + 392715392632*a(n-45) - 1340492097519*a(n-46) + 232821401640*a(n-47) - 826705375498*a(n-48) + 647431897720*a(n-49) + 135749654673*a(n-50) + 219795393014*a(n-51) - 234573750004*a(n-52) + 418157523937*a(n-53) + 41414581168*a(n-54) + 125583442065*a(n-55) - 6913751181*a(n-56) - 129465671085*a(n-57) - 20110973917*a(n-58) - 78196719452*a(n-59) + 4138743738*a(n-60) + 13082546835*a(n-61) + 6714667341*a(n-62) + 4221522441*a(n-63) + 2618944650*a(n-64) - 725257314*a(n-65) + 643911849*a(n-66) - 276231222*a(n-67) + 66784419*a(n-68) - 62178597*a(n-69) - 14348907*a(n-70)
\end{quote}
As of the ``Last modified'' line on the live entry (Nov 12 2015 22:58:50, revision 4) this is still
recorded as empirical, and nothing on the entry records it as proved.

\section{The object is a matching}

The grid has $W=4$ rows across and holds each of $0,1,\dots,LW-1$ exactly once. The value
$v$ has a HOME cell, namely its place in the plain row-major filling, and putting $v$ in cell
$(i,j)$ changes its index by $(i-h_i,\,j-h_j)$ where $(h_i,h_j)$ is that home cell. The entry
allows the index changes
\[
(i-h_i,\,j-h_j)\;\in\;\{(-2,-1),\ (-1,-1),\ (-1,2),\ (0,1),\ (1,0),\ (2,2)\},
\]
the list being the entry's own, read with the sign convention its wording uses. An admissible
arrangement is therefore exactly a perfect matching in the bipartite graph joining each cell to
the values it may hold.

\section{The matching is a walk}

Let $R=2$ be the largest $|i-h_i|$ on the list. A value whose home row is $h$ can only be
placed in rows $h-R,\dots,h+R$. So once the cells of rows $0,\dots,i$ are filled, every
value with home row at most $i-R$ has been used up, and the only home rows that can be
partly used are $i-R+1,\dots,i+R$ --- exactly $2R$ of them. Take as the state the tuple of
$2R$ bitmasks recording which values of those home rows are already placed. One step fills
the cells of one further row, drawing each of its $W$ values from a home row within $R$
and at an allowed column offset, and the home row that closes at that step must come out
full. Nonexistent home rows --- before the first and after the last --- are carried as
full, which makes the starting and accepting states the same one.

Hence $a$ is a walk count on $S=65536$ vertices,
\[
a(n)\;=\;\iota^{\!\top}M^{\,n+1}\tau ,
\]
the exponent fixed by the entry's own shape and checked against its published terms; in
particular $a$ is $C$-finite.

\section{The proof}

\begin{lemma}\label{lem:crit}
Let $q(t)=t^{70}-\sum_i c_it^{70-i}$ be the characteristic polynomial of the
conjectured recurrence. The residual $a(n)-\sum_ic_ia(n-i)$ equals
$\iota^{\!\top}M^{\,j}q(M)\tau$ for the corresponding walk index $j$, so the recurrence
holds from some point on if and only if $\iota^{\!\top}M^{j}q(M)\tau=0$ for all large $j$;
and it is enough to examine $j<S$.
\end{lemma}

\begin{proof}
Factor $M^{j}$ out of $M^{j+70}-\sum_ic_iM^{j+70-i}$. For the bound, $M^{S}$
is an integer combination of $I,M,\dots,M^{S-1}$ by Cayley--Hamilton, so the numbers
$u_j=\iota^{\!\top}M^{j}q(M)\tau$ obey the monic recurrence given by the characteristic
polynomial of $M$; $S$ consecutive zeros force every later one, and the last nonzero $u_j$
pins the threshold exactly.
\end{proof}

\begin{theorem}
\[
a(n)\;=\;7\,a(n-1) - 3\,a(n-2) + 82\,a(n-3) + 283\,a(n-4) - 1154\,a(n-5) + 4042\,a(n-6) + 17289\,a(n-7) - 22185\,a(n-8) + 32767\,a(n-9) + 188026\,a(n-10) + 232472\,a(n-11) - 37483\,a(n-12) + 1140203\,a(n-13) - 150763\,a(n-14) + 7561993\,a(n-15) - 22181996\,a(n-16) - 58336867\,a(n-17) - 54942280\,a(n-18) - 61937490\,a(n-19) + 338713215\,a(n-20) - 616216060\,a(n-21) - 750478274\,a(n-22) + 3004008536\,a(n-23) + 1358005083\,a(n-24) - 3290415548\,a(n-25) - 5332099311\,a(n-26) + 18835203759\,a(n-27) - 7620056042\,a(n-28) + 24271726336\,a(n-29) - 21686365629\,a(n-30) + 23070567658\,a(n-31) - 557733773\,a(n-32) + 145268540287\,a(n-33) + 48269697816\,a(n-34) + 145549126524\,a(n-35) + 209871190816\,a(n-36) + 171564703835\,a(n-37) + 322482128177\,a(n-38) + 392962918650\,a(n-39) + 322521808333\,a(n-40) + 509779938296\,a(n-41) - 612220254822\,a(n-42) + 615594893007\,a(n-43) - 1246199796257\,a(n-44) + 392715392632\,a(n-45) - 1340492097519\,a(n-46) + 232821401640\,a(n-47) - 826705375498\,a(n-48) + 647431897720\,a(n-49) + 135749654673\,a(n-50) + 219795393014\,a(n-51) - 234573750004\,a(n-52) + 418157523937\,a(n-53) + 41414581168\,a(n-54) + 125583442065\,a(n-55) - 6913751181\,a(n-56) - 129465671085\,a(n-57) - 20110973917\,a(n-58) - 78196719452\,a(n-59) + 4138743738\,a(n-60) + 13082546835\,a(n-61) + 6714667341\,a(n-62) + 4221522441\,a(n-63) + 2618944650\,a(n-64) - 725257314\,a(n-65) + 643911849\,a(n-66) - 276231222\,a(n-67) + 66784419\,a(n-68) - 62178597\,a(n-69) - 14348907\,a(n-70)
\]
for every $n>70$.
\end{theorem}

\begin{proof}
The vector $w=q(M)\tau$ was formed by $70$ matrix--vector products in exact integer
arithmetic and $\iota^{\!\top}M^{j}w$ evaluated for $j=0,1,\dots$, again exactly; those
integers vanish from the index corresponding to $n=70+1$ onwards and the run continued
until the working vector was identically zero, settling every larger $j$ at once.
\end{proof}

\section{Verification}

Three checks, each independent of the algebra above.

First, the reading of the entry's wording was tested against the entries themselves by brute
force before any of this was built, and separately for each of the three sign conventions the
family uses: ``$(\pm,\pm)$ $a,b$'' signs each coordinate independently, ``$\pm(.,.)$ $a,b$''
signs the PAIR and may write the second coordinate negative, and ``directed $a,b$'' is taken
literally. Enumerating all $720$ arrangements of a $2\times3$ grid reproduces the first term
of an entry of each convention. Testing only one of them would have proved only that one:
reading the second convention as the first silently drops displacements, and thirty-eight
entries disagreed with their own published data until the conventions were separated.

Second, the transfer model reproduces all $20$ terms this entry publishes, exactly, in
integer arithmetic.

Third, the annihilation test of Lemma~\ref{lem:crit} was carried out in exact integer
arithmetic on the whole vector rather than on a sampled prefix.

\begin{thebibliography}{9}
\bibitem{oeis} The OEIS Foundation, \emph{The On-Line Encyclopedia of Integer
Sequences}, \url{https://oeis.org}, 2026. Sequence A264418.
\bibitem{stanley} R.~P.~Stanley, \emph{Enumerative Combinatorics}, Volume 1, 2nd ed.,
Cambridge University Press, 2011. (Transfer-matrix method, Section 4.7.)
\bibitem{lp} L.~Lov\'asz and M.~D.~Plummer, \emph{Matching Theory}, AMS Chelsea, 2009.
\bibitem{horn} R.~A.~Horn and C.~R.~Johnson, \emph{Matrix Analysis}, 2nd ed., Cambridge
University Press, 2013.
\end{thebibliography}

\end{document}
